\begin{annexitem}{annex:p13-johnson-mccall-immediate}{thm:p13-johnson-mccall-immediate}{The native immediate inferences and their dependencies}
\begin{proof}
All transformations in this proof are classical tautology instances plus
modus ponens, uniform term substitution, declared definition interchange,
and categorical double-negation interchange. In particular, from asserted
\(p\to q\) and \(q\to r\), the tautology
\((p\to q)\to((q\to r)\to(p\to r))\) gives implication composition.
Contraposition uses the tautology \((p\to q)\to(\neg q\to\neg p)\).
These do not require modal necessitation.

First A2 and A13 give \(Iaa\). A4 with \(c=b\) gives
\(Abb\to(Iba\to Iab)\); A1 gives \(Abb\), hence \(Iba\to Iab\)
and its uniformly exchanged converse. A4 with its original \(a=b\)
gives \(Abc\to(Ibb\to Ibc)\); use \(Ibb\) to get assertoric
subalternation, then rename to \(Aab\to Iab\).

For modal oppositions, expand \(M\) and the categorical \(E,O\).
For example \(MOab=\neg L\neg\neg Aab\), equivalent by categorical
double-negation interchange to \(\neg LAab\). Similarly
\(MIab=\neg LEab\), \(MEab=\neg LIab\), and
\(MAab=\neg LOab\). The other four stated equivalences follow by
Boolean double negation and reversal of these biconditionals.

Necessary I-conversion is A11. A6 with \(b=a\), followed by \(Aaa\),
gives \(LEca\to LEac\); rename the terms to obtain necessary E-conversion.
Contrapose the two necessary conversions, use the just-established
oppositions, and exchange the terms to obtain possible E- and I-conversion.

A7 with \(b=a,c=b\) gives \(LAab\to(Iaa\to LIab)\).
The already derived \(Iaa\) gives \(LAab\to LIab\).
A8 with \(b=a,c=b\) gives \(LEab\to(Iaa\to LOab)\), hence
\(LEab\to LOab\). Contraposition with the opposition equivalences
gives \(MEab\to MOab\) from the former and \(MAab\to MIab\)
from the latter. This repairs the source's p.~267 A11 citation without
assuming the desired LA-to-LI result.

Compose LA-to-LI with I-conversion to get \(LAab\to LIba\).
Compose E-conversion with LE-to-LO to get \(LEab\to LOba\).
Contrapose these two implications and exchange \(a,b\), obtaining
\(MEab\to MOba\) and \(MAab\to MIba\), respectively.

Necessary A/I/O subordination is A12/A13/A14. From LO-to-O by
contraposition and opposition get \(Aab\to MAab\), so A1 gives
\(MAaa\). Reverse A8's two-antecedent conditional to obtain
\(MAac\to(Iab\to MIbc)\): negated LO is MA and negated LE is MI.
Put \(c=a\), use \(MAaa\), and compose with possible I-conversion.
This gives \(Iab\to MIab\). Contraposition gives \(LEab\to Eab\).
Finally contrapose LI-to-I and LA-to-A to get \(Eab\to MEab\) and
\(Oab\to MOab\). Thus all eight subordinations have actual dependencies.
\end{proof}
\end{annexitem}

\begin{annexitem}{annex:p13-johnson-mccall-rejections}{thm:p13-johnson-mccall-rejections}{The ten formula-correct rejection deductions}
\begin{proof}
We use only the three displayed rejected formulas and local assertion
dependencies. A classical equivalence between two nested conditionals
provides an asserted implication in each direction; rejection detachment
therefore transfers rejection between them. This justifies antecedent
interchange and either contrapositive form used below. If \(q\to s\)
is asserted, the classical tautology
\((q\to s)\to((p\to(s\to r))\to(p\to(q\to r)))\)
and detachment justify rejection antecedent weakening. If \(s\to r\)
is asserted, the corresponding implication from the conditional ending
in \(s\) to the one ending in \(r\) justifies rejection consequent
strengthening. Oppositions are those already proved in the body result.

\emph{R2 route.} Put \(t=Aac\to(LAca\to LAac)\).
The implication \(t\to\sigma_2\) is a propositional tautology: if \(t\)
holds, adding the antecedents in \(\sigma_2\) cannot defeat its consequence.
From rejected \(\sigma_2\), rejection detachment gives rejected \(t\).
A5 with \(b=c,c=a\), then antecedent interchange, asserts
\(u=Aac\to(LAca\to LAaa)\). By A0,
\(u\to(v\to t)\) is asserted, where
\(v=Aac\to(LAaa\to LAac)\). Detach \(u\), then apply rejection
detachment with rejected \(t\), obtaining rejected \(v\).
It is the \(b=a\) instance of Barbara XLL
\(Abc\to(LAab\to LAac)\), so rejection substitution gives that general
formula. Reversing its inner implication gives
\(Abc\to(MOac\to MOab)\); exchange \(b,c\) to get Baroco XMM.
The other contrapositive form is
\(MOac\to(LAab\to Obc)\); exchanging \(a,b\) yields Bocardo MLX.

\emph{R3 route.} Put \(t=LAcc\to(MAac\to Aac)\).
The implication \(t\to\sigma_3\) is A0, not a Barbara application.
Rejected \(\sigma_3\) therefore gives rejected \(t\). This is the
\(b=c\) instance of Barbara LMX, so rejection substitution rejects
\(LAbc\to(MAab\to Aac)\). Inner reversal gives
\(LAbc\to(Oac\to LOab)\); exchange \(b,c\) for Baroco LXL.
The other reversal gives \(Oac\to(MAab\to MObc)\); exchange \(a,b\)
for Bocardo XMM.

\emph{R4 route.} Put \(t=MAbc\to(LAab\to Iac)\).
The implication \(t\to\sigma_4\) is A0, so rejected \(\sigma_4\)
gives rejected \(t\) by rejection \emph{detachment}, not rejection
substitution. This is Barbari MLX. Interchange antecedents to get rejected
\(LAab\to(MAbc\to Iac)\). Since \(Ica\to Iac\) is asserted,
consequent strengthening rejects \(LAab\to(MAbc\to Ica)\).
Exchange \(a,c\) to obtain Bramantip LMX.
Reversing the original \(t\) gives
\(Eac\to(LAab\to LObc)\); exchange \(a,b\) to get Felapton XLL.
Assertoric \(Ebc\to Obc\) follows by contraposing A-to-I.
Antecedent weakening now rejects \(Obc\to(LAba\to LOac)\),
which is Bocardo XLL. It is not the distinct Baroco XLL formula.

All ten listed formulas have thus been rejected with the stated basis.
No claim that the partial R1--R4 list rejects every nonthesis is made.
\end{proof}
\end{annexitem}

\begin{annexitem}{annex:p13-johnson-mccall-models}{thm:p13-johnson-mccall-models}{Checking the native structures and counterformulas}
\begin{proof}
The rows in each displayed model partition its domain and every n-plus set
is nonempty, proving B1 and B2. We check the universally quantified remaining
conditions by the distinct row types; the unmentioned terms duplicate the
a-row, so add no new type.

In \(\mathcal M_1\), all n-minus sets are empty. B3's antecedent cannot
hold because actual positive sets are nonempty; the same is true of B5's
inclusion antecedent. B6's overlap antecedent is false. For B4, the only
term whose n-plus region can contain an actual positive set is b, with
n-plus equal to the entire domain. Every n-plus(x) is nonempty and therefore
overlaps n-plus(b), as required.

In \(\mathcal M_2\), the distinct actual positive sets are \(\{1\}\)
and \(\{2,3\}\). Neither is included in an n-minus region, whose only
nonempty value is \(\{2\}\), so B3 and B5 are vacuous. For B4 the
inclusion antecedent requires y and z of a/b type; any x overlapping
\(\{1\}\) is also of that type. Their n-plus values all equal \(\{1\}\).
For B6 the inclusion antecedent again requires z and y of a/b type,
and their n-minus values are equal. Its conclusion therefore repeats
the overlap already assumed.

In \(\mathcal M_3\), the distinct actual positive sets are \(\{1,2\}\)
and \(\{3\}\). Neither is included in an n-minus region, whose only
nonempty value is \(\{2\}\); B3 and B5 are vacuous. B4's inclusion
requires y and z of b/c type, and any actual positive set overlapping
that type is itself of that type. Their n-plus values equal \(\{3\}\).
B6's inclusion also requires y and z of b/c type; their n-minus regions
are equal, so the conclusion repeats the assumed overlap.

In \(\mathcal M_4\), each actual positive set contains 1, whereas no
n-minus region contains 1. B3 and B5 are therefore vacuous. Each actual
positive set has a contingent positive member not in any n-plus region
that could contain it: the actual sets are \(W,\{1,3,4\},\{1,4\}\),
and n-plus sets are \(\{1,2\},\{1,3\},\{1\}\). No inclusion from
an actual positive set into any n-plus region holds. Thus B4 and B6
are vacuous too. This checks all six base conditions in all four structures.

For Barbara XLL in \(\mathcal M_1\), \(Abc\) is true because both
actual b and c equal W; \(LAab\) is true because \(\{1,2\}\subseteq W\);
\(LAac\) is false because \(\{1,2\}\nsubseteq\{3\}\).
For Celarent XLL in \(\mathcal M_2\), actual b and c are disjoint,
so \(Ebc\) is true; \(LAab\) is true since \(\{1\}\subseteq\{1\}\);
\(LEac\) is false since \(\{1\}\nsubseteq\varnothing\).
For Camestres LXL in \(\mathcal M_3\), \(LAcb\) is true since
\(\{3\}\subseteq\{3\}\); \(Eab\) is true since
\(\{1,2\}\cap\{3\}=\varnothing\); \(LEac\) is false since
\(\{1,2\}\nsubseteq\{2\}\).
For Baroco XLL in \(\mathcal M_4\), \(Acb\) is true since
\(\{1,4\}\subseteq\{1,3,4\}\); \(LOab\) is true because
\(\{1,2\}\cap\{2\}=\{2\}\); \(LOac\) is false because
\(\{1,2\}\cap\{3\}=\varnothing\).

For a row permutation, assign each target term the four regions of its
chosen source row. Every universal base-condition instance becomes the
same condition on those source rows; it was already checked. Extend
unmentioned target terms by the target a-row. Thus the permuted structure
is a model, including for repeated row types.
Swap a and b in \(\mathcal M_1\). Then \(MEbc\) is true because
n-plus(b) and n-plus(c) are disjoint, \(LIab\) is true because
n-plus(a)=W overlaps n-plus(b), and \(Oac\) is false because both
actual a and c equal W. This is Ferio MLX.
For Baroco XMM, take target a from source a, target b from source c,
and target c from source b. Both actual b and c are W, so \(Acb\) is
true. \(MOab=\neg LAab\) is true since \(\{1,2\}\nsubseteq\{3\}\);
\(MOac=\neg LAac\) is false since \(\{1,2\}\subseteq W\).
Every counterformula is consequently false under ordinary material
implication in a model satisfying the native conditions.
\end{proof}
\end{annexitem}

\begin{annexitem}{annex:p13-johnson-q-collapse}{thm:p13-johnson-q-collapse}{The two native conditional-collapse deductions}
\begin{proof}
We work in the unmodified Q-L-X-M assertion basis, where A20,A21,A28
are still present and A2 is ordinary Iaa. The already supplied native
immediate-inference proof applies with Iaa available directly.
Implication composition, antecedent interchange, and reversal below
are tautology instances with modus ponens. No unrestricted Q-distribution
or necessitation is assumed.

We first justify the required Q conversions. A17 with its first two
subject variables identified and A1 yields \(QAab\to QEab\).
A22 gives the reverse implication. A18 with the same identification,
ordinary Iaa, and modus ponens gives \(QAab\to QIab\).
A25 then gives \(QAab\to QIba\). A23,A24 give complementary
equivalence of QI and QO. These supply every Q operation used below.

For the first collapse, A20 uniformly substituted gives
\(Eca\to(QAbc\to MEba)\). Necessary conversion per accidens gives
\(LAab\to LIba\); contraposition and native modal oppositions give
\(MEba\to MOab\). Composition yields
\(Eca\to(QAbc\to MOab)\). Interchange and inner reversal give
\(QAbc\to(LAab\to Ica)\), since negated MOab is LAab and
negated Eca is Ica. A7 uniformly substituted gives
\(LAab\to(Ica\to LIcb)\). The tautology that composes these
two conditionals under their shared antecedents therefore gives
\(QAbc\to(LAab\to LIcb)\).
On the other hand QA-to-QE, A26, and possible E-conversion give
\(QAbc\to QEbc\to MEbc\to MEcb\), understood as successive
asserted implications. Native opposition identifies MEcb with
\(\neg LIcb\). Thus \(QAbc\to\neg LIcb\) is asserted.
For any formula r, the classical tautology
\[
 (p\to(q\to s))\to((p\to\neg s)\to(p\to(q\to r)))
\]
with p=QAbc, q=LAab, s=LIcb, and two modus-ponens steps proves
\(QAbc\to(LAab\to r)\).

For the second collapse, A21 uniformly substituted gives
\(Eac\to(QIba\to MObc)\). The earlier QA-to-QIba implication
strengthens its second antecedent to give
\(Eac\to(QAab\to MObc)\). Reversal and interchange yield
\(LAbc\to(QAab\to Iac)\). Compose with assertoric I-conversion
to get \(LAbc\to(QAab\to Ica)\).
A18 uniformly substituted gives \(QAab\to(Ica\to QIcb)\).
Propositional composition under the two antecedents yields
\(LAbc\to(QAab\to QIcb)\). Compose its conclusion with A25,
A23, and A28 to get, successively, QIbc, QObc, and MObc.
The last is \(\neg LAbc\) by native opposition.
Thus \(LAbc\to(QAab\to\neg LAbc)\) is asserted.
The tautology
\((p\to(q\to\neg p))\to(p\to(q\to r))\), followed by
modus ponens, proves \(LAbc\to(QAab\to r)\).
Both deductions have the stated arbitrary consequent, while retaining
their two antecedents. Neither deduction eliminates those antecedents
to make r itself an unconditional theorem.
\end{proof}
\end{annexitem}

\begin{annexitem}{annex:p13-johnson-q-nonempty}{thm:p13-johnson-q-nonempty}{The omitted nonemptiness and the essential-existence boundary}
\begin{proof}
For the literal printed-condition counterexample, take W with one object,
and for every term put n-minus=W and the other three regions empty.
Then actual positive and q are both empty. The regions partition W.
BQ2 has empty sets included in n-minus on both sides. BQ3/BQ4
have false overlap antecedents; BQ5 has false n-plus overlap; BQ6
has empty inclusion on all three positions; BQ7 has a false overlap
antecedent; and BQ8 has a disjoint empty n-plus region. Hence all
printed BQ conditions hold. Nevertheless Iaa is false because the
actual positive region is empty. Nonempty W alone does not prevent this.

For the completed-model example take W with one object, and for every
term put q-plus=W and the other three regions empty. Actual positive
and q both equal W, so ordinary term nonemptiness holds. BQ2 has a
false inclusion antecedent, as do BQ3 and BQ4; BQ5 also has a false
inclusion antecedent. BQ6 is inclusion of W in W. BQ7's concluding
overlap is true. BQ8 has disjoint empty n-plus and n-minus regions.
Thus this is a completed Q-model. Iaa is true, but LIaa is false
because n-plus(a) is empty. The completion repairs ordinary existence
without assuming essential positive instances.
\end{proof}
\end{annexitem}

\begin{annexitem}{annex:p13-johnson-q-soundness}{thm:p13-johnson-q-soundness}{All retained axioms, transformations, and the A29 addition}
\begin{proof}
Fix an arbitrary completed Q-model. All semantic inclusions and overlaps
below are in its one object domain. To prove a two-antecedent conditional,
suppose its two antecedents true and show its consequent true. This is
ordinary material implication, with the source's p.~272 consequent-value
slip corrected as explicitly stated in the body.

A0 is valid because Boolean tautologies remain true for every valuation
of their well-formed categorical and modal propositional units.
A1 is reflexive inclusion of actual positive regions. A2 is true because
each actual positive region is nonempty. A3 is transitivity of inclusion.
For A4, an object in actual b and actual a belongs also to actual c
when actual b is included in actual c.

For A5, actual a is included in actual b, which is included in n-plus(c),
so actual a is included in n-plus(c). For A6, actual c included in
n-minus(b) gives actual b included in n-minus(c) by BQ2; actual a
included in actual b is therefore included in n-minus(c).
A7 is exactly BQ3 with y=b,z=c,x=a. A8 is BQ4 with the same
assignment. A9 is BQ5 with z=c,y=b,x=a. For A10 choose a witness
in n-plus(b) intersect n-minus(c). Actual b is included in n-plus(a);
since n-plus(b) is included in actual b, the same witness lies in
n-plus(a) intersect n-minus(c), proving LOac.
A11 follows from symmetry of intersection. A12 follows because
n-plus(b) is included in actual b. For A13 the common essential-positive
witness is a common actual-positive witness. For A14 a witness in
n-plus(a) intersect n-minus(b) is an actual a outside actual b,
so Aab is false and Oab true.

A15 is BQ6 with y=b,z=c,x=a. For A16 the QIab antecedent gives
actual a overlapping q(b) or q(a) overlapping actual b; with actual b
included in q(c), BQ7 yields QIac. For A17 actual a included in
actual b, which is included in q(c), gives actual a included in q(c),
the truth condition for QEac. For A18 an actual a/b witness, together
with actual b included in q(c), gives actual a overlapping q(c),
so QIac. For A19, actual b included in actual c and actual a included
in q(b) imply n-plus(a) disjoint from n-minus(c) by BQ8.
Thus LOac is false and MAac true. A20 and A21 are deleted and
are not being smuggled into this proof.

A22 is valid because QA and QE have the same declared truth condition.
A23 and A24 are valid because QI and QO have the same truth condition.
For A25 the two disjuncts in that condition exchange under swapping a,b,
using symmetry of overlap, so QI conversion is valid.
For A26, actual a included in q(b) implies n-plus(a) disjoint from
n-plus(b), since q(b) is disjoint from n-plus(b). LIab is false,
and MEab true. For A27 argue contrapositively. If MIab were false,
actual a would be included in n-minus(b). BQ2 would make actual b
included in n-minus(a). The partition would then exclude both actual a
overlapping q(b) and q(a) overlapping actual b. QIab would be false.
This proves QIab-to-MIab. A28 is deleted and has not been used.

For A29, suppose QAbc and MIab true. The latter is precisely the
failure of actual a inclusion in n-minus(b), so choose an actual a
outside n-minus(b). By the partition for b this object lies either
in actual b or in q-minus(b). In the first case actual a overlaps
actual b; actual b included in q(c) gives actual a overlapping q(c),
so QIac. In the second case actual a overlaps q(b); actual b included
in q(c) and the retained printed BQ7 give actual a overlapping q(c),
again QIac. The cases exhaust the witness's possible regions.
Thus A29 is valid without replacing MIab by Iab.

For uniform term substitution, compose each of the four region functions
with the substitution map on terms. Every base condition and ordinary
nonemptiness remain true because they were universal over term assignments.
The truth of the original formula in this composed model equals the truth
of its substituted instance in the original model, by induction on the
restricted formula grammar. Thus substitution preserves validity, including
when terms are identified. Modus ponens preserves truth by Boolean
implication. Definition interchange preserves truth by the declared
definitional clauses and their categorical normalization; induction on a
well-formed propositional context carries that identity to the whole formula.
Categorical double-negation interchange under L or Q preserves truth by the
explicit parity-normalization clause, and outside modalities by Boolean
double negation. These are all assertion transformations; there are no
contingency rejection transformations to validate.

Finally induct on the finite length of an assertion deduction. Its axiom
lines are valid by the preceding cases, and every permitted transformation
preserves validity. Hence every asserted formula is valid in every completed
Q-model. The argument includes all retained axioms and A29, rather than only
the illustrative axiom checks printed in the source. It proves no converse.
\end{proof}
\end{annexitem}

\begin{annexitem}{annex:p13-johnson-t3-comparison}{thm:p13-johnson-t3-comparison}{Inclusion, strictness, and the equivalent negative condition}
\begin{proof}
A T3-model already satisfies J-model B1 and B2. For B3, assume
actual z included in n-minus(y) and actual x included in actual y.
BT4 gives actual y included in n-minus(z), hence so is actual x.
For B4, assume actual y included in n-plus(z) and actual x overlapping
actual y. Applying BT3 with x and y exchanged gives actual y overlapping
n-plus(x). Its witness also lies in n-plus(z), giving the required
n-plus(x)/n-plus(z) overlap. For B5 use the same exchanged BT3 witness,
but actual y is now included in n-minus(z). This yields overlap with
n-minus(z), not the source proof's slipped n-plus(z).
For B6, actual z included in n-plus(y) makes n-minus(y) included in
n-minus(z) by BT5; any assumed n-plus(x)/n-minus(y) witness therefore
is the required n-plus(x)/n-minus(z) witness. All J-conditions follow.

The already validated J-model \(\mathcal M_1\) has actual a=\(\{1,2\}\),
actual c=W, and n-plus(c)=\(\{3\}\). Actual a overlaps actual c but
does not overlap n-plus(c). BT3 fails. Thus inclusion is strict, with
an actual counterexample rather than the erroneous M4 example.

For the last claim, BT4 implies B3 by the first argument of this proof.
Conversely take universally quantified B3, suppose actual z included
in n-minus(y), and instantiate its x parameter by y. Actual y is
included in itself, so B3 gives actual y included in n-minus(z).
That is exactly BT4. No additional existence or model-class assumption
is required for this equivalence of the two universally quantified conditions.
\end{proof}
\end{annexitem}

\begin{annexitem}{annex:p13-johnson-q-countercases}{thm:p13-johnson-q-countercases}{The three source Q-models and five failures}
\begin{proof}
All displayed rows partition their domain and have nonempty actual positive
regions. Duplication of the a-row for other terms adds no new type.
We explicitly check the remaining Q-conditions.

For M5 the actual regions are a=\(\{1,2\}\), b=c=\(\{1,3\}\).
None is included in any n-minus region, so BQ2 and BQ4 have false
inclusion antecedents. The only n-plus region containing an actual region
is n-plus(c)=\(\{1,3\}\), containing actual b and c. Every n-plus(x)
overlaps that region, proving BQ3. None of the n-plus regions overlaps
n-minus(c)=\(\{2\}\), so BQ5 is vacuous. The only actual-to-q inclusion
is actual a included in q(b)=\(\{1,2\}\). No actual region is included
in q(a)=\(\{2\}\), so BQ6 is vacuous. Every actual region overlaps
q(b) at 1 or 2, proving BQ7 for its only possible y=a,z=b instance.
For BQ8 its x must be a and y must be b; actual b is included only
in actual b/c. n-plus(a)=\(\{1\}\) is disjoint from both n-minus(b)
and n-minus(c), as required.

For M6 the actual regions are a=\(\{1,2\}\), b=\(\{1,3\}\), c=\(\{4\}\).
The necessary-negative inclusions are exactly actual a/b into n-minus(c)
and actual c into n-minus(a/b), so their reversals prove BQ2.
The only actual-to-n-plus inclusion is c into itself. Only c overlaps
actual c, and n-plus(c) overlaps itself, proving BQ3.
For BQ4, actual a/b included in n-minus(c) can overlap only actual
a/b; their essential-positive witnesses 1/3 belong to n-minus(c).
Actual c included in n-minus(a/b) can overlap only c; its witness 4
belongs to both those negative regions. This exhausts BQ4. BQ5's
inclusion has z=y=c, so its concluding overlap repeats its assumed
overlap. The only actual-to-q inclusion is a into q(b)=\(\{1,2\}\);
no actual region is included in q(a)=\(\{2\}\), proving BQ6 vacuously.
For BQ7 with y=a,z=b, the premise overlap can involve x=a/b, and
each of their actual regions overlaps q(b); x=c supplies neither
premise overlap. For BQ8 x=a,y=b and z=b; n-plus(a) is disjoint
from n-minus(b)=\(\{4\}\). All conditions hold.

For M7, a/c have actual region \(\{1\}\), n-plus \(\{1\}\),
n-minus \(\{3\}\), and q \(\{2\}\); b has actual region \(\{2,3\}\),
n-plus \(\{2\}\), n-minus empty, and q \(\{1,3\}\).
No actual region is included in an n-minus region; BQ2 and BQ4
are vacuous. Actual-to-n-plus inclusions involve only a/c types, and
only those types overlap their actual region, proving BQ3. No n-plus
region overlaps n-minus(a/c), so BQ5 is vacuous. Actual-to-q inclusions
have y=a/c,z=b. No actual region is included in q(a/c), proving BQ6
vacuously. The BQ7 premise can then involve only x=b, whose actual region
overlaps q(b) at 3, proving the conclusion. For BQ8 x=a/c,y=b,z=b;
n-minus(b) is empty. Thus M7 is a completed Q-model too.

In M5 LAbc is true since actual b equals n-plus(c). QAab is true
since actual a equals q(b). QAac is false since q(c) is empty.
This falsifies Barbara LQQ. Also QObc is true because q(b) overlaps
actual c at 1; MObc is false because LAbc is true. Thus the deleted
QO-to-MO axiom has a genuine counterinstance.
In M6 LEac and QAab are true, while QObc is false: actual b cannot
overlap empty q(c), and q(b)=\(\{1,2\}\) misses actual c=\(\{4\}\).
This is Felapton LQQ with minor b and major c.
In M7 QAcb and LAac are true, while Iab is false. Hence
\(QAcb\to(LAac\to Iab)\), an instance of the first arbitrary-consequent
collapse, is false. Finally QEab is true because actual a is included
in q(b), whereas QEba is false because \(\{2,3\}\nsubseteq\{2\}\).
This is a counterinstance to ordinary universal contingent conversion,
without disputing complementary QA/QE equivalence.
\end{proof}
\end{annexitem}

\begin{annexitem}{annex:p13-johnson-q-generated}{thm:p13-johnson-q-generated}{The thirteen A29-based assertion deductions}
\begin{proof}
QA-to-QE follows from A17 with identified subject variables and A1;
QE-to-QA is A22. QA-to-QI follows from A18 with identified subjects
and ordinary Iaa. QI-to-QO and its reverse are A23/A24; QI-conversion
is A25. Possible I-conversion and MA-to-MI were fully derived in
the native immediate-inference proof and remain derivable with Iaa
available directly. All implication compositions and antecedent changes
below use classical tautologies and modus ponens, not modal distribution.

Start with A29, Darii QMQ:
\(QAbc\to(MIab\to QIac)\).
Replace its first antecedent by the stronger QEbc using QE-to-QA,
and its consequent by QOac using QI-to-QO, obtaining Ferio QMQ.
Replace MIab by MIba using possible I-conversion to obtain Datisi QMQ,
\(QAbc\to(MIba\to QIac)\).
Replace MIba by MAba using MA-to-MI to obtain Darapti QMQ.
Apply QE-to-QA and QI-to-QO to these last two to obtain respectively
Ferison QMQ and Felapton QMQ.
From the original A29 replace MIab by MAab to obtain Barbari QMQ;
apply QE-to-QA and QI-to-QO to obtain Celaront QMQ. These are
the eight stated QMQ assertions with their exact premise placements.

In Datisi QMQ exchange a,c uniformly. The result is
\(QAba\to(MIbc\to QIca)\). Interchange antecedents and convert
the QI conclusion to obtain Disamis MQQ,
\(MIbc\to(QAba\to QIac)\).
Replace MIbc by MAbc to obtain Darapti MQQ.
In the original Darii QMQ exchange a,c uniformly to get
\(QAba\to(MIcb\to QIca)\); interchange antecedents and convert
the conclusion, obtaining Dimaris MQQ,
\(MIcb\to(QAba\to QIac)\).
Replace MIcb by MAcb to obtain Bramantip MQQ.
In this last formula replace QAba by QEba and QIac by QOac,
obtaining Camenop MQQ, \(MAcb\to(QEba\to QOac)\).
These are the five stated MQQ assertions. None used deleted A20,
A21, or A28, nor a rejection rule under Q.
\end{proof}
\end{annexitem}
